\section{Method and Setup}
\label{sec:method}

\paragraph{Corpus.}
We screened \NCandidates{} candidate programs and kept \NPrograms{}:
\NTestSuite{} from the LLVM test suite, \NPolyBench{} from PolyBench, \NCBench{} from
cBench~\cite{milepost}, and \NOwn{} small hand-written kernels. A program is kept only if it
compiles, terminates within the time window, passes UndefinedBehaviorSanitizer and
AddressSanitizer, and produces deterministic output. The output is recorded as a reference digest.
Multi-file programs are linked at the IR level. cBench programs run on their first dataset with a
loop wrapper tuned to about 0.5\,s. Of the kept programs, \NRuntimeReliable{} have run times long
and stable enough for wall-clock timing.

\paragraph{Passes and interaction score.}
We use \NPasses{} LLVM passes: \PassList{}. This gives \NPairs{} ordered pairs of distinct passes.
For a program state $P$ and metric $m$ (lower is better), let
$\Delta(B \mid P) = (m(P) - m(B(P)))/m(P)$ be the gain of running pass $B$ on $P$. The
interaction of the ordered pair $(A,B)$ is
\[ I(A,B,P) = \frac{m(A(P)) - m(B(A(P)))}{m(P)} - \Delta(B \mid P), \]
that is, how much running $A$ first changes $B$'s gain. Positive values mean $A$ enables $B$;
negative values mean $A$ disables it. Pairs with $|I| \le \Threshold$ are \emph{weak}. The search
objective and the interaction metric for all hypotheses is the IR instruction count, with no-op
intrinsics (\texttt{llvm.lifetime.*}, \texttt{llvm.dbg.*}) excluded.

\paragraph{Snapshots.}
Interactions are measured at \NSnapshots{} program states. S0 is raw IR from
\texttt{clang -O1 -Xclang -disable-llvm-passes}. This form has no \texttt{optnone} and no
compiler-added \texttt{noinline}, so the inliner can act. S1--S4 are prefixes of LLVM's own
printed \texttt{default<O2>} pipeline: after early simplification (S1), after the CGSCC
inliner and function simplification (S2), before the late module passes (S3), and the full
pipeline (S4). In total there are \NExperiments{} experiments (\NExpFailed{} failed), yielding
\NInteractionRows{} interaction values over all metrics.

\paragraph{Machine-code metrics.}
Every IR variant is compiled by one fixed backend (\texttt{clang -O2} with IR passes disabled).
We measure \texttt{.text} size for x86-64, and for AArch64 from compile-only objects. We also
measure executed instructions with Valgrind's cachegrind, and wall-clock time. Wall-clock time is
the median of 20 runs after one warm-up, pinned to one core, run serially and alone, in a randomised
order. Every measured binary's output is checked against the reference: \OutputFailures{}
mismatches. \texttt{-Os} and \Oz{} are applied with clang's \texttt{optsize}/\texttt{minsize}
attributes, so they differ from \Otwo{}.

\paragraph{Static graphs and predictors.}
All static graphs and models for a program are built from the \emph{other} programs only
(leave-one-program-out, LOPO). Hyperparameters are tuned on a validation split of the training
programs. We use three static graphs:
\begin{itemize}
\item the \emph{static graph}: the mean interaction per ordered pair, plus the mean single-pass gain;
\item a \emph{DCO-style} reimplementation: positive, negative and neutral pair classes, with sub-sequences along positive edges searched by OpenTuner;
\item a \emph{Pan-style} reimplementation: a synergy graph whose weighted random walks seed and mutate a genetic algorithm.
\end{itemize}
These are reimplementations in the style of the originals, not the original systems. The
predictor is a gradient-boosted tree model of the residual from the static graph. Its inputs are
the pass pair and \NReps{} representations of the IR the passes run on: hand-crafted counts,
LLVM function properties, an opcode histogram, and IR2Vec~\cite{ir2vec}. The snapshot name is not
an input, so the model applies to any IR state.

\paragraph{Ordering under a compile budget.}
The budget unit \E is one \emph{evaluation}: a not-yet-seen candidate sequence is applied to the program's S0 IR in one \texttt{opt} invocation and its
objective is measured. A candidate sequence is drawn from the \NPasses{} passes; its length is
capped at 16 passes, and each pass may appear at most twice. We also count all \texttt{opt} invocations, including feature extraction,
and wall-clock time. Each method records its best-so-far result at
$\E \in \{1,5,10,25,50,100,200\}$. The methods are:
\begin{itemize}
\item graph-based: \emph{static graph}, \emph{predict-once} (the LOPO model's prediction of the program's S0 graph),
and \emph{oracle graph} (the program's own measured graph; this needs \NPairs{} extra measurements and is a
headroom reference, not a usable method), all using the same capped beam builder;
\item \emph{re-prediction}: the model is re-applied to the features of every intermediate IR state;
\item \emph{graph-guided} search: graph proposals verified by measurement;
\item search: Pan-style, DCO-style, OpenTuner~\cite{opentuner}, a genetic algorithm (GA) and random search, each with 5 seeds;
\item \emph{measured greedy}, as a reference.
\end{itemize}

\paragraph{Pre-registration and statistics.}
Before the pilot and before any full result, we froze the hypotheses, tests, thresholds and
analysis choices in a pre-registration document (SHA-256 \texttt{\PreregShaShort\ldots}, full hash
in the artifact). There are four primary hypotheses. Each is a one-sided Wilcoxon signed-rank test~\cite{wilcoxon1945}
over programs, with Holm correction~\cite{holm1979} across exactly these four:
\begin{itemize}
\item H1: oracle graph vs.\ static graph at $\E=10$, on per-program log ratios; supported if Holm $p<0.05$ and geomean ratio $\le 0.98$.
\item H2: per-program Spearman $\rho$ between S0 and S2 interactions over pairs non-weak in either; supported if $\rho<0.5$ and median $\le 0.5$.
\item H3: LOPO model vs.\ static graph, per-program MAE; supported if residual $R^2 \ge 0.10$.
\item H4: re-prediction vs.\ the strongest other method at $\E=10$, same thresholds as H1.
\end{itemize}
All other results are exploratory. They are reported with bootstrap 95\% confidence intervals and
raw, uncorrected $p$-values.

\paragraph{Deviations from the pre-registration.}
One. After the pilot, a system library (zlib) was installed so that two cBench programs could be
screened. Both then failed the sanitizer screen, so the kept corpus is unchanged. Only the
manifest file's drop reasons, and therefore its hash, changed. Two operational changes did not
alter any analysis choice: fewer parallel model fits to stay within memory, and a restart of the
timing stage after a process check misfired.
